\chapter{Couverture et inventaire en pratique}\label{ch:hf:hedging-and-inventory-in-practice}

À quatre heures moins dix, un teneur de marché est long d'un panier de quarante titres qu'il n'a jamais choisis. Ses clients ont vendu dans un après-midi en baisse, ses cotations se sont infléchies
pour compenser chaque position, comme le prévoient les modèles des chapitres 3 et 4, et pourtant le portefeuille détient encore 1.6~million de dollars d'actions, dont 830\,000~\$ en exposition équivalente à l'indice.
Il peut vendre les actions, vendre contre elles le contrat à terme sur indice, ou les conserver pendant la nuit et repartir le lendemain en position longue. Dans
le marché simulé de ce chapitre, vendre les actions coûte 440~\$~; vendre le contrat à terme coûte 83~\$ et laisse un écart-type nocturne de 3\,100~\$~;
garder le portefeuille expose à une perte nocturne de 14\,800~\$ dans un cas sur cent. Pendant la journée, la même question se pose toutes les quelques secondes, et y répondre
à chaque mouvement coûte 14~\% de la fourchette que gagne le desk.

\section{Choisir l'instrument de couverture}

L'inventaire d'un teneur de marché est un sous-produit~: il détient ce que les clients lui ont vendu. Les modèles d'inventaire des chapitres 3 et 4 le gèrent par les cotations,
qui ne coûtent rien mais agissent lentement. Une couverture agit aussitôt mais traverse une fourchette.

\begin{definition}[Instrument de couverture]\label{def:hf:hedging-and-inventory-in-practice:instrument}
Un \emph{instrument de couverture}\index{instrument de couverture} est un titre qu'un teneur de marché échange, non pour gagner une fourchette mais pour compenser le risque des positions
qu'il a acquises en faisant du market making~; il est choisi pour le coût de sa négociation et pour la part du risque des positions qu'il supprime.
\end{definition}

Le choix porte sur des instruments qui suppriment des parties différentes du risque à des prix différents. Dans un modèle à un facteur, le titre $i$ a pour prix $S_i$, un
bêta $\beta_i$ par rapport au facteur de marché et une volatilité quotidienne idiosyncratique $\sigma_i$~; le facteur a une volatilité quotidienne $\sigma_F$ et le contrat à terme sur indice
reproduit celui-ci.

\begin{definition}[Inventaire en équivalent delta]\label{def:hf:hedging-and-inventory-in-practice:delta}
L'\emph{inventaire en équivalent delta}\index{inventaire en équivalent delta} d'un portefeuille est son exposition exprimée dans un instrument de référence~: chaque
position multipliée par sa sensibilité à la référence, sommée sur les instruments. Face à un indice, un portefeuille d'actions $q_i$ a
\[
D=\sum_i q_i\,S_i\,\beta_i
\]
en dollars d'exposition à l'indice~; un ETF compte avec son propre bêta, une option avec son delta multiplié par le bêta de son sous-jacent.
\end{definition}

La variance quotidienne du portefeuille est $D^2\sigma_F^2+\sum_i (q_iS_i\sigma_i)^2$. Vendre $D$ dollars de contrat à terme supprime le premier terme à un coût $c_F|D|$,
où $c_F$ est le coût du contrat à terme par dollar négocié~; en ramenant chaque action à zéro, on supprime les deux termes à un coût $c_S\sum_i|q_iS_i|$. Avec une charge $\lambda$
sur chaque dollar carré de variance détenu pendant un horizon $h$, le choix compare trois nombres (\cref{lst:hf:hedging-and-inventory-in-practice:choose})~:
\[
\lambda h\,\bigl(D^2\sigma_F^2+\textstyle\sum_i(q_iS_i\sigma_i)^2\bigr),\qquad c_F|D|+\lambda h\textstyle\sum_i(q_iS_i\sigma_i)^2,\qquad c_S\textstyle\sum_i|q_iS_i|.
\]
Le contrat à terme l'emporte lorsque le portefeuille est concentré sur le facteur et de grande taille, les actions lorsque le risque idiosyncratique est fort ou l'horizon long, et ne
rien faire lorsque le portefeuille est petit. Un contrat à terme comporte aussi des risques que le modèle à un facteur omet~: sa base par rapport à l'indice au comptant, les bêtas du portefeuille estimés avec
une erreur d'estimation, et, pour un portefeuille d'actions hors indice, une exposition aux facteurs que le contrat à terme ne reproduit pas. Le chapitre~9 du livre~9 traite le
ratio de couverture estimé.

\section{Couverture partielle et déclencheurs de couverture}

La couverture n'est que rarement une affaire de tout ou rien. Si le seul coût est proportionnel, le problème à une période admet une forme fermée.

\begin{proposition}[La zone morte]\label{prop:hf:hedging-and-inventory-in-practice:dead}
Vendre $x$ dollars de contrat à terme contre une exposition $D$, à un coût $c_F|x|$ et avec une charge de risque $\lambda\sigma_F^2(D-x)^2$, est optimal pour
\[
x^\ast=\operatorname{sign}(D)\,\max\Bigl(|D|-\frac{c_F}{2\lambda\sigma_F^2},\,0\Bigr):
\]
le teneur de marché se couvre jusqu'au bord d'une zone morte de demi-largeur $c_F/(2\lambda\sigma_F^2)$, et non à l'intérieur.
\end{proposition}

\begin{proof}
Pour $0<x<D$, la dérivée de $c_Fx+\lambda\sigma_F^2(D-x)^2$ vaut $c_F-2\lambda\sigma_F^2(D-x)$, nulle en $D-x=c_F/(2\lambda\sigma_F^2)$~; l'objectif est
convexe, et aucune couverture au-delà de $D$ ou de signe opposé n'améliore l'objectif. Si $|D|$ est à l'intérieur de la demi-largeur, la dérivée en $x=0$ est déjà positive ou nulle.
\end{proof}

Au cours d'une journée, l'exposition évolue en continu à mesure que les exécutions arrivent, et la réponse à une période devient une bande~: ne rien faire tant que l'exposition non couverte reste
dans $\pm b$, revenir au bord de la bande quand elle en sort. Le chapitre~10 du livre~4 dérive la bande optimale pour une exposition qui diffuse
avec une volatilité $\sigma_D$, une charge de risque courante $\gamma X^2$ et un coût proportionnel $c$~: $b^\ast=\bigl(3c\sigma_D^2/(4\gamma)\bigr)^{1/3}$. Ici
$\gamma=\lambda\sigma_F^2$. La racine cubique rend la bande robuste~: doubler le coût l'élargit de 26~\%.

Un déclencheur de couverture peut aussi être un événement plutôt qu'un niveau~: une grosse exécution, une exécution venant d'un client réputé informé, ou un mouvement du marché qui
fait bondir le risque de l'exposition. La discipline est la même~: chaque couverture doit rapporter plus en risque supprimé qu'elle ne coûte.

\section{Inventaire à travers des portefeuilles corrélés}

Une firme qui fait du market making sur les actions de l'indice, sur ses ETF et sur ses contrats à terme détient une seule exposition au facteur, répartie entre plusieurs portefeuilles. Compenser les
\omterm{def:hf:hedging-and-inventory-in-practice:delta}{inventaires en équivalent delta} avant de couvrir est la couverture la moins chère qui soit~: une position longue dans un portefeuille compense une position courte dans un autre, et la firme ne couvre que la
somme. Le desk qui fait du market making sur les ETF (chapitre~12) peut souvent se couvrir avec ses propres inventaires d'actions~; le delta du desk d'options (chapitre~19) s'ajoute
à ce même nombre.

La compensation repose sur deux conditions. Les portefeuilles doivent être valorisés et sommés quasiment en temps réel, ce qui est un problème d'ingénierie~; et les risques résiduels
que la compensation laisse (idiosyncratiques, de base, de fuseau horaire) doivent rester dans les limites propres à chaque portefeuille.

Les données empiriques indiquent que l'inventaire pèse sur les prix et sur la liquidité. Sur douze ans de données quotidiennes des spécialistes du New York Stock Exchange, de 1994 à
2005, Hendershott et Menkveld ont mesuré des pressions sur les prix liées à l'inventaire des intermédiaires, de 49 points de base en moyenne, avec une demi-vie de 0.92
jour~: 17 points de base et 0.54 jour dans les plus grandes valeurs, 118 points de base et 2.11 jours dans les plus petites. Comerton-Forde, Hendershott, Jones,
Moulton et Seasholes ont constaté, sur onze ans de données de ces mêmes spécialistes, que les fourchettes à l'échelle du marché et à celle des titres s'élargissent lorsque les spécialistes détiennent
de grandes positions ou perdent de l'argent, le plus fortement lorsque les inventaires sont importants ou les résultats médiocres. Et le teneur de marché à haute fréquence étudié par Menkveld en 2013 a
gagné un bénéfice sur la fourchette et perdu sur ses positions~: sur chaque transaction, 1.55~\euro de gain sur la fourchette net de frais, dont 0.68~\euro sont partis dans le positionnement.

\section{Fin de journée et nuit}

\begin{definition}[Mise à plat de fin de journée]\label{def:hf:hedging-and-inventory-in-practice:flatten}
La \emph{mise à plat de fin de journée}\index{mise à plat de fin de journée} (end-of-day flattening) est la réduction des positions d'un teneur de marché vers zéro avant la clôture, par ses
cotations, par des ordres passés sur le marché continu ou dans l'enchère de clôture, ou en couvrant ce qui reste avec un instrument négocié la nuit.
\end{definition}

La nuit change le problème~: les actions ne peuvent pas être négociées avant l'ouverture, l'écart d'ouverture n'est pas une diffusion mais un saut à queues épaisses, et
le portefeuille n'a aucun revenu de fourchette à lui opposer. Quatre façons de mettre à plat, de la moins chère à la plus chère~:

\begin{itemize}
\item \emph{Par les cotations.} Renforcer le skew dans la dernière heure. Dans le tutoriel, des cotations inclinées quatre fois plus fort à partir de 15:00 réduisent de moitié le portefeuille brut à
la clôture, de 1.32~million de dollars à 670\,000~\$, tandis que le desk réalise le même nombre d'exécutions.
\item \emph{Dans l'enchère de clôture.} Aucune fourchette à franchir, un prix unique~; le risque tient au déséquilibre propre à l'enchère, et un gros ordre fait bouger la clôture
(livre~1, chapitre~13).
\item \emph{Sur le marché continu} avant la clôture~: la demi-fourchette plus l'impact.
\item \emph{Garder les actions et vendre le contrat à terme pendant la nuit}~: le risque de facteur disparaît, les écarts idiosyncratiques demeurent.
\end{itemize}

\begin{omfigure}
\begin{tikzpicture}
\begin{axis}[width=11cm,height=6.4cm,xlabel={\small risque intrajournalier (milliers de dollars par jour)},ylabel={\small coût de couverture (milliers de dollars par jour)},
  xmin=0,xmax=10,ymin=-1,ymax=14.5,tick label style={font=\footnotesize},grid=major,grid style={gray!25},table/col sep=comma]
\addplot[omThm,only marks,mark=*,mark size=1.8pt] table[x=risk,y=cost]{figdata/hft/11-hedging-and-inventory-in-practice/frontier.csv};
\begin{scope}[font=\scriptsize]
\node[anchor=west] at (axis cs:2.7,12.76) {actions};
\node[anchor=west] at (axis cs:4.6,6.70) {future à chaque mouvement};
\node[anchor=east] at (axis cs:4.35,3.14) {bande 100~k\$};
\node[anchor=west] at (axis cs:4.9,2.3) {bande 250~k\$};
\node[anchor=north] at (axis cs:5.44,0.6) {bande 500~k\$};
\node[anchor=south] at (axis cs:6.99,0.6) {bande 1~M\$};
\node[anchor=south,align=center] at (axis cs:8.9,0.3) {aucune\\couverture};
\end{scope}
\end{axis}
\end{tikzpicture}

\omcaption{Coût de couverture face au risque intrajournalier (écart type du P\&L à la minute, ramené à une journée) pour sept règles~: pas de couverture, le future maintenu
à moins d'un demi-contrat de l'\omterm{def:hf:hedging-and-inventory-in-practice:delta}{inventaire en équivalent delta} du carnet, le future hors des bandes de non-transaction de 100\,000~\$ à 1~million de dollars, et chaque titre
ramené à 20\,000~\$ dans son propre titre. Quarante titres, dix journées simulées, revenu du spread de 46\,800~\$ par jour. Données~:
\texttt{hf\_hedging.frontier}.}\label{fig:hf:hedging-and-inventory-in-practice:frontier}
\end{omfigure}

\section{Tutoriel~: quarante titres, un facteur, un future}\label{tut:hf:hedging-and-inventory-in-practice:tut}

\begin{tutorial}
\textbf{Objectif.} Mesurer le coût de chaque règle de couverture et le risque qu'elle laisse sur une journée simulée d'exécutions portant sur quarante titres, puis prendre la décision «~dix contre quatre~».
\textbf{État final~:} les \cref{fig:hf:hedging-and-inventory-in-practice:frontier,fig:hf:hedging-and-inventory-in-practice:overnight} et le tableau
de la décision «~dix contre quatre~».
\begin{tutsteps}
\item \textbf{La journée.} \texttt{firm.mmhedge.FlowDay} tire quarante titres (cours de 20 à 150~\$, bêtas de 0.6 à 1.4, volatilité idiosyncratique de 1 à 2~\% par jour,
facteur à 1~\%), 0.05 exécution par seconde et par titre, portant sur 100 actions, à un demi-spread d'un cent~; les cotations sont biaisées en fonction de l'inventaire, et des clients vendent après une baisse du
facteur (une inclinaison de 0.15 par écart type du mouvement de la dernière minute). Cette inclinaison fait évoluer ensemble les inventaires des titres~: l'exposition du
carnet en équivalent delta présente une volatilité quotidienne de 2.5~millions de dollars. Simuler les exécutions plutôt que quarante carnets d'ordres complets maintient la journée à une fraction
de seconde~; le harnais du chapitre~1 est l'outil à utiliser lorsque les exécutions dépendent du carnet.
\item \textbf{Le choix} de l'\omterm{def:hf:hedging-and-inventory-in-practice:instrument}{instrument de couverture} (\cref{lst:hf:hedging-and-inventory-in-practice:choose}).
\omcode{code/firm/mmhedge/firm_mmhedge.py}{47}{56}{Coût plus la charge sur la variance restante, pour les trois choix.}\label{lst:hf:hedging-and-inventory-in-practice:choose}
\item \textbf{Les règles} sur le future~: à moins d'un demi-contrat (50\,000~\$ de notionnel), ou une bande de non-transaction.
\omcode{code/firm/mmhedge/firm_mmhedge.py}{125}{144}{La position en future le long de l'exposition de la journée, et son coût.}\label{lst:hf:hedging-and-inventory-in-practice:band}
\item \textbf{La frontière.} Dix journées, sept règles, future à un demi-point de base par dollar, actions à 1.7 points de base (\cref{fig:hf:hedging-and-inventory-in-practice:frontier}).
\item \textbf{Dix contre quatre.} Prendre le carnet à 15~h~50 le premier jour et comparer la mise à plat sur les dix dernières minutes (demi-spread plus impact
en racine carrée), la vente du future pour la nuit (aller-retour~: un point de base) et son maintien, avec des écarts de Student-$t$ à quatre degrés de liberté (facteur 0.6~\%,
titres 1~\%).
\end{tutsteps}
\textbf{Pour aller plus loin.} Ajouter une base entre le future et le facteur~; estimer les bêtas sur un mois de rendements simulés au lieu
d'utiliser les vrais~; regrouper dans $D$ un second carnet (un ETF).
\end{tutorial}

La frontière (\cref{fig:hf:hedging-and-inventory-in-practice:frontier})~: sans couverture, le risque intrajournalier s'élève à 8\,540~\$ par jour. Maintenir le future à moins d'un demi-contrat
le ramène à 4\,410~\$, mais génère 2\,682 contrats échangés et coûte 6\,700~\$ par jour, soit 14~\% du revenu du spread. Une bande de 250\,000~\$ laisse
4\,760~\$ de risque pour 1\,880~\$. Les actions éliminent aussi le risque idiosyncratique (2\,525~\$) pour 12\,760~\$. Avec une charge de $\lambda=10^{-4}$ par
dollar au carré, la bande de 250\,000~\$ obtient le meilleur score (42\,570~\$ de revenu du spread, moins les coûts et la charge, contre 39\,640~\$ sans couverture et 38\,260~\$ pour le future à
chaque mouvement), et la bande donnée par la formule vaut 284\,000~\$.

\begin{omfigure}
\begin{tikzpicture}
\begin{axis}[width=11cm,height=5cm,xlabel={\small perte overnight (milliers de \$)},ylabel={\small probabilité},xmin=-25,xmax=35,ymin=0,
  yticklabel style={/pgf/number format/fixed,/pgf/number format/precision=2},scaled y ticks=false,
  tick label style={font=\footnotesize},grid=major,grid style={gray!25},table/col sep=comma,
  legend style={font=\scriptsize,at={(0.98,0.95)},anchor=north east}]
\addplot[omThm,thick,const plot mark mid] table[x=loss,y=keep]{figdata/hft/11-hedging-and-inventory-in-practice/overnight.csv};
\addplot[omDef,thick,dashed,const plot mark mid] table[x=loss,y=future]{figdata/hft/11-hedging-and-inventory-in-practice/overnight.csv};
\draw[omThm,dotted] (axis cs:14.78,0) -- (axis cs:14.78,0.06);
\draw[omDef,dotted] (axis cs:7.36,0) -- (axis cs:7.36,0.06);
\legend{conservé,{conservé, future vendu}}
\end{axis}
\end{tikzpicture}

\omcaption{Perte overnight du carnet détenu à 15~h~50 (1.57~million de dollars en brut, 834\,000~\$ en équivalent delta), conservé tel quel ou avec 834\,000~\$ de future
vendus, par classes de 1\,000~\$~; écarts de Student-$t$ à quatre degrés de liberté, facteur 0.6~\% et titres 1~\% par nuit~; 200\,000 tirages. Pointillés~: les
quantiles à 99~\%, 14\,780~\$ et 7\,360~\$. Données~: \texttt{hf\_hedging.ten\_to\_four}.}\label{fig:hf:hedging-and-inventory-in-practice:overnight}
\end{omfigure}

\begin{center}
\small
\setlength{\tabcolsep}{4pt}%
\begin{tabular}{@{}>{\raggedright\arraybackslash}p{4.6cm}>{\centering\arraybackslash}p{2.3cm}>{\centering\arraybackslash}p{2.3cm}>{\centering\arraybackslash}p{2.3cm}>{\centering\arraybackslash}p{2.3cm}@{}}
\toprule
à 15~h~50 & coût & écart type overnight & perte à 99~\% & expected shortfall à 99~\%\\
\midrule
mise à plat sur les dix dernières minutes & 438~\$ & 0 & 0 & 0\\
vente du future pour la nuit & 83~\$ & 3\,086~\$ & 7\,358~\$ & 8\,813~\$\\
conserver & 0 & 5\,879~\$ & 14\,776~\$ & 19\,795~\$\\
\bottomrule
\end{tabular}
\end{center}

Avec la même charge, les trois choix obtiennent 438~\$, 1\,035~\$ et 3\,454~\$~: la mise à plat l'emporte, car les écarts idiosyncratiques que le future ne peut pas éliminer
sont comptés à la variance d'une nuit, sans spread pour les compenser. Incliner les cotations à partir de 15~h~00 aurait réduit de moitié la facture de mise à plat (163~\$
contre 381~\$ en moyenne sur les dix jours).

\section{Construction~: le module de couverture}\label{bld:hf:hedging-and-inventory-in-practice:mmhedge}

\begin{build}
\bfield{Rôle} Exprimer les positions d'un teneur de marché comme une seule exposition à un facteur, choisir ce que l'on couvre, et appliquer les règles de bande, de déclenchement et de fin de journée.
\bfield{Interface} \texttt{delta\_equivalent}, \texttt{book\_var}, \texttt{choose(q, p, beta, sf, se, c\_fut, c\_stk, lam, horizon)},
\texttt{partial\_hedge}, \texttt{optimal\_band} (depuis \texttt{firm.impulse}), \texttt{FlowDay(n, seed)} avec
\texttt{inventory(stock\_band, c\_stk, eod\_from, eod\_skew)}, \texttt{hedge\_future(D, rule, band, notional, c)}, \texttt{day\_pnl},
\texttt{overnight}, \texttt{flatten\_cost}.
\bfield{Règles} Coûts par dollar échangé~; le future s'échange par contrats entiers~; la marge de la bande, et non zéro, est la cible après un déclenchement~; les écarts overnight
sont à queues épaisses.
\bfield{Tests d'acceptation} \texttt{code/firm/mmhedge/tests/}~: exposition et variance à la main, et le choix dans trois cas construits~; la zone morte est
l'argmin numérique~; les règles sur le future maintiennent l'exposition dans leurs bornes, et la bande échange un cinquième autant~; l'inclinaison fait évoluer les titres ensemble, et la
bande sur les actions plafonne chaque titre~; l'écart type overnight concorde avec la formule~; coût de mise à plat calculé à la main~; la bande concorde avec la forme fermée~; un
biais plus fort à partir de 15~h~00 réduit de moitié le carnet à la clôture.
\bfield{Pour aller plus loin} Erreur d'estimation de la base et des bêtas~; plusieurs carnets regroupés~; une mise à plat à la clôture avec un risque de déséquilibre.
\end{build}

\begin{omsources}
\item T.~Hendershott, A.~J.~Menkveld, Price pressures, \emph{Journal of Financial Economics} 114(3), 2014, 405--423.
\item C.~Comerton-Forde, T.~Hendershott, C.~M.~Jones, P.~C.~Moulton, M.~S.~Seasholes, Time variation in liquidity: the role of market-maker
inventories and revenues, \emph{Journal of Finance} 65(1), 2010, 295--331.
\item A.~J.~Menkveld, High frequency trading and the new market makers, \emph{Journal of Financial Markets} 16(4), 2013, 712--740.
\end{omsources}

\section{Exercices}

\begin{exercise}[$\star$]\label{exo:hf:hedging-and-inventory-in-practice:1}
Un carnet est long de 1\,000 actions d'un titre à 50~\$ de bêta 1.2 et court de 500 actions d'un titre à 100~\$ de bêta 0.8. Calculer son \omterm{def:hf:hedging-and-inventory-in-practice:delta}{inventaire en équivalent delta},
et déterminer combien de futures de 50\,000~\$ la règle du demi-contrat vend.
\end{exercise}

\begin{exercise}[$\star$]\label{exo:hf:hedging-and-inventory-in-practice:2}
Avec $c_F=0.5$ point de base, $\lambda=10^{-4}$ et $\sigma_F=1~\%$, calculer la demi-largeur de la zone morte et déterminer quelle part d'une exposition de 833\,747~\$ est couverte.
\end{exercise}

\begin{exercise}[$\star$]\label{exo:hf:hedging-and-inventory-in-practice:3}
Calculer le coût de la vente de 833\,747~\$ de future à 15~h~50 puis de son rachat à l'ouverture, à 0.5 point de base dans chaque sens.
\end{exercise}

\begin{exercise}[$\star\star$]\label{exo:hf:hedging-and-inventory-in-practice:4}
Le future supprime le risque de facteur. Pourquoi le carnet couvert présente-t-il encore un écart type overnight de 3\,086~\$~?
\end{exercise}

\begin{exercise}[$\star\star$]\label{exo:hf:hedging-and-inventory-in-practice:5}
Calculer la bande optimale pour $\lambda=10^{-5}$, $3\times10^{-5}$ et $10^{-4}$ avec $\sigma_D=\$2.476$ million par jour, et expliquer pourquoi la bande
se resserre lorsque l'aversion au risque augmente.
\end{exercise}

\begin{exercise}[$\star\star$]\label{exo:hf:hedging-and-inventory-in-practice:6}
Renforcer l'asymétrie des cotations à partir de 15~h~00 réduit de moitié le carnet à la clôture. Pourquoi cette opération est-elle presque gratuite, et dans quels cas ne le serait-elle pas~?
\end{exercise}

\begin{exercise}[$\star\star\star$]\label{exo:hf:hedging-and-inventory-in-practice:7}
\emph{Programmation.} Doubler le coût du future (\texttt{frontier(1e-4)}). Quelle bande simulée obtient le meilleur score, et que dit la formule~?
\end{exercise}

\begin{exercise}[$\star\star\star$]\label{exo:hf:hedging-and-inventory-in-practice:8}
\emph{Trouver la faille.} «~Nous vendons un future chaque fois que le bêta du carnet varie d'un contrat. Cela minimise notre risque, donc c'est nécessairement la bonne politique.~»
\end{exercise}

\section{Problème~: Quarante titres à dix heures moins dix}

\begin{problem}[{Problème du week-end --- quarante titres à dix heures moins dix}]\label{pb:hf:hedging-and-inventory-in-practice:1}
Un teneur de marché sur actions termine sa journée long d'un panier qu'il n'a pas choisi et doit décider comment se couvrir au cours de la journée et ce qu'il garde pendant la nuit.

\medskip\noindent\textbf{Partie I --- L'exposition.}
\begin{enumerate}
\item Définir un \omterm{def:hf:hedging-and-inventory-in-practice:instrument}{instrument de couverture} et les coûts qui le choisissent.
\item Définir l'\omterm{def:hf:hedging-and-inventory-in-practice:delta}{inventaire en équivalent delta} d'un carnet d'actions par rapport à un indice.
\item Écrire la variance quotidienne du carnet dans le modèle à un facteur.
\item Pourquoi les clients qui vendent sur un marché en baisse font-ils évoluer ensemble les titres d'un teneur de marché~?
\end{enumerate}

\medskip\noindent\textbf{Partie II --- Règles.}
\begin{enumerate}[resume]
\item Énoncer et démontrer la zone morte (\cref{prop:hf:hedging-and-inventory-in-practice:dead}).
\item Donner la bande optimale et sa valeur pour $\lambda=10^{-4}$.
\item Nommer trois risques d'une couverture par future que le modèle à un facteur omet.
\item Pourquoi compenser les positions entre carnets avant de couvrir, et à quelles conditions~?
\end{enumerate}

\medskip\noindent\textbf{Partie III --- Mesures.}
\begin{enumerate}[resume]
\item Donner le coût et le risque intrajournalier des quatre règles principales.
\item Quelle règle obtient le meilleur score pour $\lambda=10^{-4}$, et à quel point la formule est-elle proche~?
\item Qu'ont montré les études empiriques sur l'inventaire et les prix, et sur l'inventaire et les spreads~?
\item Combien le teneur de marché de Menkveld a-t-il gagné sur le spread et perdu sur ses positions~?
\end{enumerate}

\medskip\noindent\textbf{Partie IV --- Le verdict.}
\begin{enumerate}[resume]
\item Énoncer le \emph{résultat nommé}~: le coût et le risque résiduel des trois choix à 15~h~50, et la distribution des pertes overnight de la conservation du carnet.
\item Quel choix l'emporte à coût égal, et pourquoi~?
\item Définir la \omterm{def:hf:hedging-and-inventory-in-practice:flatten}{mise à plat de fin de journée} (flattening) et énumérer ses quatre voies.
\item Que change l'asymétrie à partir de 15~h~00 dans la facture de mise à plat~?
\item Quand conserver le carnet pendant la nuit serait-il justifié~?
\item Comment la réponse change-t-elle si les bêtas sont estimés avec erreur~?
\item Qu'ajouterait un carnet d'options à l'\omterm{def:hf:hedging-and-inventory-in-practice:delta}{inventaire en équivalent delta}~?
\item En une phrase~: quand une couverture vaut-elle son coût~?
\end{enumerate}
\end{problem}

\section{Questions d'entretien}

\begin{interviewq}[$\star$ \iqroles{trader}]\label{iq:hf:hedging-and-inventory-in-practice:1}
Vous êtes long de 5~millions de dollars d'actions bancaires à 15~h~55. Quelles sont vos options, et laquelle retenez-vous~?
\end{interviewq}

\begin{interviewq}[$\star\star$ \iqroles{chercheur}]\label{iq:hf:hedging-and-inventory-in-practice:2}
Comment estimeriez-vous les bêtas servant à couvrir un carnet de quarante titres avec le future sur indice, et comment testeriez-vous la couverture~?
\end{interviewq}

\begin{interviewq}[$\star\star$ \iqroles{chercheur}]\label{iq:hf:hedging-and-inventory-in-practice:3}
Pourquoi un coût proportionnel produit-il une bande sans transaction, et un coût fixe un retour brutal à zéro~?
\end{interviewq}

\begin{interviewq}[$\star\star$ \iqroles{risque}]\label{iq:hf:hedging-and-inventory-in-practice:4}
Trois desks signalent chacun une petite exposition à l'indice, dans les limites. Que pourrait-il mal tourner au niveau de la firme~?
\end{interviewq}

\begin{interviewq}[$\star\star$ \iqroles{développeur}]\label{iq:hf:hedging-and-inventory-in-practice:5}
Comment concevriez-vous le service qui calcule en temps réel l'\omterm{def:hf:hedging-and-inventory-in-practice:delta}{inventaire en équivalent delta} de la firme sur les actions, les ETF, les futures et les options~?
\end{interviewq}

\begin{interviewq}[$\star\star\star$ \iqroles{chercheur}]\label{iq:hf:hedging-and-inventory-in-practice:6}
Comment dériveriez-vous la demi-largeur de la zone morte, et comment varie-t-elle avec l'horizon de la couverture~?
\end{interviewq}
